\section{Etapa 4 --- La red inal\'ambrica ESP-NOW}
\label{sec:et4}
\periodo{24 de mayo --- 3 de junio de 2026 \quad$\cdot$\quad 7 d\'ias de trabajo, 18 \emph{commits}}

\subsection{El problema que resuelve}

Un sism\'ografo multicanal convencional cablea todos los ge\'ofonos a un
digitalizador \'unico, y entonces la sincronizaci\'on entre canales es gratis:
hay un solo reloj. Al pasar a nodos inal\'ambricos independientes, cada nodo
tiene su propio reloj y su propio arranque, y \emph{la sincronizaci\'on pasa a
ser el problema central del sistema}.

Y no es un problema menor de comodidad. En MASW, un desfasaje entre canales se
traduce directamente en un error de velocidad de fase, que se propaga a la
curva de dispersi\'on y de ah\'i al perfil $V_S$. Un arreglo desincronizado no
da un resultado ruidoso: da un resultado \emph{equivocado y cre\'ible}.

\subsection{Arquitectura adoptada}

\begin{center}\small
\begin{tabularx}{\textwidth}{@{}l X@{}}
\toprule
\textbf{Elemento} & \textbf{Rol} \\
\midrule
Maestro ESP32 & Coordina el arreglo, sincroniza el disparo, concentra los vol\'umenes de datos y expone la interfaz hacia la PC \\
Esclavos ESP32 & Uno por nodo. Comandan al PSoC local, reciben las muestras y las retienen hasta el volcado \\
ESP-NOW & Enlace de control y sincronizaci\'on: sin asociaci\'on previa, latencia baja y predecible ---lo que hace falta para disparar, no para transferir \\
PSoC & Adquisici\'on real. El ESP no toca la se\~nal: pide, espera y transporta \\
\bottomrule
\end{tabularx}
\end{center}

\descartado{\textbf{Enlace SPI PSoC\,$\leftrightarrow$\,ESP.} La primera
versi\'on (24 de mayo) usaba SPI con el PSoC como esclavo. Se reemplaz\'o por
UART en el refactor del 31 de mayo. En la placa de agosto la subida vuelve a
cambiar, ahora a I\textsuperscript{2}C (Etapa~11): es la interfaz que
m\'as veces se revis\'o del proyecto.}

\subsection{La secuencia de sincronizaci\'on: c\'omo se lleg\'o a PRESTART}

Vale la pena registrar la evoluci\'on completa, porque es el mejor ejemplo de
c\'omo el sistema aprendi\'o a no confiar en el tiempo.

\textbf{Intento 1 --- espera fija.} \verb|ARM| $\to$ esperar $N$ ms $\to$
\verb|START|. Falla cuando ESP-NOW tarda m\'as que $N$: el esclavo recibe
\verb|START| antes de estar en \verb|ARMED| y directamente no arranca. El
s\'intoma es traicionero, porque el resto del arreglo s\'i arranca.

\textbf{Intento 2 --- \emph{probe}/ACK con medici\'on de RTT} (25 de mayo). Ya
no se supone la latencia: se mide. Mejora, pero sigue habiendo ventana entre
la confirmaci\'on y el disparo.

\textbf{Soluci\'on --- PRESTART con HOT\_WAIT} (26 de mayo). Se agrega un estado
intermedio en el que el maestro \emph{verifica} que todos los esclavos est\'an
calientes antes de emitir el disparo real:

\begin{center}\small
\texttt{IDLE $\to$ ARMING $\to$ ARMED $\to$ \textbf{PRESTART} $\to$ RUNNING $\to$ STOPPING}
\end{center}

\begin{enumerate}[itemsep=0.15em]
\item El maestro difunde \verb|CMD_ARM|.
\item Espera los \verb|ARM_ACK| con \emph{timeout}.
\item Entra en PRESTART y consulta a cada esclavo su estado (HOT\_WAIT:
consulta cada \SI{20}{\milli\second}, \emph{timeout} total
\SI{120}{\milli\second}, 3 reintentos, todo configurable).
\item Solo con todos confirmados emite el \verb|START| real.
\end{enumerate}

\seleccionado{\textbf{Y todo el mecanismo es no bloqueante.} Este es el detalle
que lo hizo funcionar: los pulsos de \emph{debug} usaban
\texttt{delayMicroseconds()}, lo que congelaba el CPU \SI{1}{\milli\second} y
hac\'ia perder callbacks de ESP-NOW ---es decir, \emph{el propio instrumento de
medici\'on causaba la p\'erdida del \emph{ACK} que se estaba tratando de
diagnosticar}. Se reescribi\'o con marcas de tiempo y verificaci\'on en el
\emph{loop}. Cambio chico en c\'odigo, decisivo en comportamiento.}

\subsection{C\'omo se midi\'o la sincronizaci\'on}

No se dio por buena: se instrument\'o. Se agreg\'o el comando
\verb|CMD_SCOPE_START| (\texttt{0x60}), un disparo \emph{falso} que hace que
todos los nodos emitan un pulso GPIO de \SI{1}{\milli\second} sin muestrear ni
responder. Con el osciloscopio en los pines de \emph{sync} de los tres nodos
(maestro GPIO\,15, esclavo\,1 GPIO\,32, esclavo\,2 GPIO\,22), la separaci\'on
entre pulsos \emph{es} la latencia de sincronizaci\'on del sistema.

Y para no medir una sola vez, el estado \verb|SCOPE_MULTI| repite el disparo
128 veces con \SI{1}{\second} de separaci\'on, de modo que el osciloscopio
acumule estad\'istica de \emph{jitter} sin intervenci\'on manual.

\dato{Objetivo de dise\~no fijado en esta etapa: todos los pulsos dentro de
\SI{1}{\milli\second} entre s\'i.}

\pendiente{El \emph{jitter} medido sobre el arreglo completo y desplegado en
campo sigue sin volcarse al documento. Ver \S\ref{sec:pendientes}.}

\subsection{Hitos de funcionamiento}

\begin{center}\small
\begin{tabularx}{\textwidth}{@{}l l X@{}}
\toprule
\textbf{Fecha} & \textbf{Mensaje} & \textbf{Qu\'e significa} \\
\midrule
26 may & \emph{``Parece que funciona la putaaaa''} & Primera transmisi\'on extremo a extremo: PSoC $\to$ ESP esclavo $\to$ ESP maestro $\to$ MATLAB \\
3 jun & \emph{``Locura parece que ya funciona super bien''} & Sistema estable tras corregir el \emph{beacon} a \SI{10}{\hertz}, \emph{store-and-forward} en \texttt{VER} y UART a 115\,200~baudios \\
\bottomrule
\end{tabularx}
\end{center}

\subsection{En paralelo: acondicionamiento del martillo (28--29 de mayo)}

La fuente de impacto necesita su propio canal, porque marca el instante cero
de la captura. Se cre\'o el proyecto PSoC \texttt{AcondicionamientoMartillo}
para el aceler\'ometro de impacto \textbf{PCB Piezotronics 086D20}, con cadena
propia: \texttt{LPF\_ref} $+$ \texttt{LPF\_ADC} $+$ referencia de
\SI{1.024}{\volt}.

Que sea un proyecto aparte no es casual: el martillo y el ge\'ofono tienen
rangos de se\~nal y anchos de banda muy distintos, y compartir cadena habr\'ia
obligado a comprometer los dos.

\subsection{Refactor de cierre (31 de mayo, PR \#2)}

El PR \#2 consolid\'o la etapa:

\begin{itemize}[itemsep=0.2em]
\item enlace PSoC\,$\leftrightarrow$\,ESP migrado de SPI a \textbf{UART};
\item contador de muestras $N$ ampliado a 16 bits;
\item m\'aquina de estados expl\'icita en el PSoC;
\item \textbf{arranque del PSoC solo por flanco f\'isico} ---no por comando---,
que es lo que garantiza que el instante cero sea el mismo en todos los nodos
independientemente de la latencia del enlace;
\item \texttt{debug\_log.h} unificado y READMEs por proyecto.
\end{itemize}

\seleccionado{\textbf{El principio que sobrevivi\'o a todo:} el disparo viaja
por software, pero el instante cero lo define un \emph{flanco f\'isico}. Todos
los rediseños posteriores ---incluida la placa de agosto, que cambia el enlace
de subida a I\textsuperscript{2}C--- mantuvieron esa separaci\'on.}

\subsection{Balance de la etapa}

\begin{center}\small
\begin{tabularx}{\textwidth}{@{}l X@{}}
\toprule
\textbf{Se qued\'o} & ESP-NOW para control y sincronizaci\'on; PRESTART/HOT\_WAIT; disparo por flanco f\'isico; instrumentaci\'on de \emph{jitter} por pulso GPIO; \texttt{SCOPE\_MULTI} para estad\'istica autom\'atica; UART hacia el PSoC; canal dedicado de martillo \\
\addlinespace[3pt]
\textbf{Se descart\'o} & Espera fija tras ARM; enlace SPI PSoC$\leftrightarrow$ESP; pulsos de \emph{debug} bloqueantes \\
\addlinespace[3pt]
\textbf{Qued\'o abierto} & La captura segu\'ia gobernada por c\'odigo C en el PSoC, con el \emph{jitter} atado a las interrupciones. Ese es el punto de partida de la Etapa~5 \\
\bottomrule
\end{tabularx}
\end{center}
